% !TEX root = ../presentation.tex
% =============================================================
%  Subsection — OHL conductors (codewords and standards)
%  Part of ac-transmission; \input from ac-transmission.tex,
%  directly after series-compensation.tex.
%
%  No figures; text and tables only.
%  Sources: ASTM B232 (ACSR), ASTM B231 (AAC), ASTM B399 (AAAC), BS 215-2
% =============================================================
\providecommand{\acfig}[1]{ac-transmission/fig/#1}

\intermezzo{OHL Conductors}{Definitions and applications}
\begin{frame}{OHL conductors --- codewords}
  \begin{columns}[T,onlytextwidth]
    \begin{column}{0.33\textwidth}
      \renewcommand{\arraystretch}{1.2}
      {\small
      \begin{tabular}{@{}l c@{}}
        \textbf{Codeword} & \textbf{Al/St strands} \\
        \hline
        Cuckoo  & 24/7  \\
        Drake   & 26/7  \\
        Mallard & 30/19 \\
        Condor  & 54/7  \\
        \hline
      \end{tabular}}\\[0.5em]
      {\color{atGray}\footnotesize The four standard strandings of
      795\,kcmil ACSR.}
    \end{column}
    \begin{column}{0.63\textwidth}
      {\small
      \begin{itemize}
        \setlength{\itemsep}{0.15em}
        \item A size alone does not define a conductor: 795\,kcmil ACSR exists
              in four standard strandings, each with a different steel content
              and so different strength, weight and sag.
        \item The codeword fixes area, stranding and core in one word.
        \item Originally Alcoa product nomenclature (1920s--30s), later
              adopted into the ASTM standards (ACSR: ASTM B232).
        \item The name family identifies the construction:
              \begin{itemize}
                \setlength{\itemsep}{0.1em}
                \item[--] ACSR / ACSS --- birds (Drake, Falcon, Rail, Cardinal)
                \item[--] AAC --- flowers (Rose, Iris, Tulip, Peony)
                \item[--] AAAC --- US city names (Akron, Amherst, Anaheim)
              \end{itemize}
        \item No size information in the name: Bluebird is 2156\,kcmil,
              Condor only 795\,kcmil.
      \end{itemize}}
    \end{column}
  \end{columns}
\end{frame}

\begin{frame}{OHL conductors --- codewords elsewhere}
  \begin{itemize}
    \setlength{\itemsep}{0.5em}
    \item BS 215 and Indian practice use the same idea with mammals:
          Dog, Wolf, Lynx, Panther, Zebra, Moose. A 400\,kV Indian line is
          specified as ``twin Moose'', the North American equivalent as
          ``bundled Drake''.
    \item Metric designs in continental Europe usually drop codewords and
          quote the aluminium and steel areas directly, e.g.\ 490/65\,mm$^2$.
    \item Practical consequence for us: a conductor must always be quoted
          with its standard as well as its codeword, since the same name can
          appear in more than one national system.
  \end{itemize}
  \vspace{0.5em}
  {\color{atGray}\footnotesize ASTM B232 (ACSR), ASTM B231 (AAC), ASTM B399 (AAAC), BS 215-2}
\end{frame}


\begin{frame}{Link collection}
  \textbf{Overhead line (OHL) conductor data}
  \begin{itemize}
    \item {\footnotesize southwire.com/wire-cable/bare-aluminum-overhead-transmission-distribution}
  \end{itemize}
  \textbf{AC cable data}
  \begin{itemize}
    \item {\footnotesize new.abb.com --- XLPE submarine cable systems (2GM5007)}
  \end{itemize}
  \textbf{PSCAD tutorials and information}
  \begin{itemize}
    \item {\footnotesize pscad.com/knowledge-base/topic-281}
    \item {\footnotesize pscad.com/knowledge-base/article/44}
  \end{itemize}
\end{frame}

\begin{frame}{Link collection}
  \textbf{Overhead line (OHL) conductor data}
  \begin{itemize}
    \item {\footnotesize southwire.com/wire-cable/bare-aluminum-overhead-transmission-distribution}
  \end{itemize}
  \textbf{AC cable data}
  \begin{itemize}
    \item {\footnotesize new.abb.com --- XLPE submarine cable systems (2GM5007)}
  \end{itemize}
  \textbf{PSCAD tutorials and information}
  \begin{itemize}
    \item {\footnotesize pscad.com/knowledge-base/topic-281}
    \item {\footnotesize pscad.com/knowledge-base/article/44}
  \end{itemize}
\end{frame}